%\newpage
%\setcounter{page}{1}
\begin{center}
\begin{figure}
\centering%
%\epsfig{file=HojaTitulo/University_of_Los_Andes_logo.xvg,scale=1}%
\includegraphics[scale=0.2]{./HojaTitulo/University_of_Los_Andes_logo.pdf}
\end{figure}
\thispagestyle{empty} \vspace*{1.5cm} \textbf{\huge
Simulaciones por Din\'{a}mica Molecular para el Estudio de Propiedades Biof\'{i}sicas de Estafiloxantina en Membranas Modelo de \textit{Staphylococcus aureus}}\\[5.0cm]
\Large\textbf{John Erick Cabrera Ramirez}\\[6.0cm]
\small Universidad de los Andes\\
Facultad de Ciencias, Departamento de F\'{i}sica\\
Bogot\'{a}, Colombia\\
%A\~{n}o\\
2020\\
\end{center}

\newpage{\pagestyle{empty}\cleardoublepage}

\newpage
\begin{center}
\thispagestyle{empty} \vspace*{-1cm} \textbf{\huge
Simulaciones por Din\'{a}mica Molecular para el Estudio de Propiedades Biof\'{i}sicas de Estafiloxantina en Membranas Modelo de \textit{Staphylococcus aureus}}\\[3.0cm]
\Large\textbf{John Erick Cabrera Ramirez}\\[3.0cm]
\small Tesis o trabajo de grado presentada(o) como requisito parcial para optar al
t\'{\i}tulo de:\\
\textbf{Mag\'{i}ster en Ciencias-F\'{i}sica}\\[2.0cm]
Director(a):\\
%T\'{\i}tulo (Ph.D., Doctor, Qu\'{\i}mico, etc.) y nombre del director(a)\\[2.0cm]
Ph. D. Chad Leidy\\[2.0cm]
L\'{\i}nea de Investigaci\'{o}n:\\
%Nombrar la l\'{\i}nea de investigaci\'{o}n en la que enmarca la tesis  o trabajo de investigaci\'{o}n\\
Biof\'{i}sica de membranas\\
%Grupo de Investigaci\'{o}n:\\
%Biof\'{i}sica de membranas\\
[2.5cm]
Universidad de los Andes\\
Facultad de Ciencias, Departamento de F\'{i}sica\\
Bogot\'{a}, Colombia\\
%A\~{n}o\\
2020\\
\end{center}

\newpage{\pagestyle{empty}\cleardoublepage}

\newpage
\thispagestyle{empty} \textbf{}\normalsize
\\\\\\%
\textbf{(Dedicatoria o un lema)}\\[4.0cm]

\begin{flushright}
\begin{minipage}{8cm}
    \noindent
        \small
        Su uso es opcional y cada autor podr\'{a} determinar la distribuci\'{o}n del texto en la p\'{a}gina, se sugiere esta presentaci\'{o}n. En ella el autor dedica su trabajo en forma especial a personas y/o entidades.\\[1.0cm]\\
        Por ejemplo:\\[1.0cm]
        A mis padres\\[1.0cm]\\
        o\\[1.0cm]
        La preocupaci\'{o}n por el hombre y su destino siempre debe ser el
        inter\'{e}s primordial de todo esfuerzo t\'{e}cnico. Nunca olvides esto
        entre tus diagramas y ecuaciones.\\\\
        Albert Einstein\\
\end{minipage}
\end{flushright}

\newpage{\pagestyle{empty}\cleardoublepage}

\newpage
\thispagestyle{empty} \textbf{}\normalsize
\\\\\\%
\textbf{\LARGE Agradecimientos}
\addcontentsline{toc}{chapter}{\numberline{}Agradecimientos}\\\\
Esta secci\'{o}n es opcional, en ella el autor agradece a las personas o instituciones que colaboraron en la realizaci\'{o}n de la tesis  o trabajo de investigaci\'{o}n. Si se incluye esta secci\'{o}n, deben aparecer los nombres completos, los cargos y su aporte al documento.\\

\newpage{\pagestyle{empty}\cleardoublepage}